\section{Finite-state diagnostics for the twisted spectral gap}
\label{sec:eta-gap}

The strict inequality
\(\eta=\max_{\varrho\neq\mathbf 1}\rho(A_\varrho)<\lambda\) is the
finite-state specialization in which the renormalized Perron-boundary package
has no non-trivial boundary phases; if \(G\) is trivial, this empty maximum is
\(0\).  This section gives the finite-state diagnostic needed for that
strict-gap branch and describes the equality case \(\eta=\lambda\).  The equality case is rigid
in the precise sense needed here: a non-trivial channel has Perron spectral
radius exactly when the cocycle admits an equivariant line field.  The resulting
loop-character diagnostic feeds \Cref{cor:strict-gap-specialization} and the older
computational formulas
\Cref{thm:class-mertens-explicit,thm:class-euler-product-mertens}.

\begin{lemma}[Equality in the Hilbert-space triangle inequality]
  \label{lem:hilbert-triangle-equality}
  Let \(x_1,\ldots,x_q\) be non-zero vectors in a Hilbert space.  If
  \[
    \left\|\sum_{a=1}^q x_a\right\|=\sum_{a=1}^q\|x_a\|,
  \]
  then there is a unit vector \(\xi\) and positive real numbers
  \(c_1,\ldots,c_q\) such that \(x_a=c_a\xi\) for every \(a\).
\end{lemma}

\begin{proof}
  Set \(s:=\sum_a x_a\).  The hypothesis and non-zeroness of the \(x_a\)'s
  imply \(s\neq0\).  For every \(a\),
  \[
    \operatorname{Re}\langle x_a,s\rangle\le \|x_a\|\,\|s\|,
  \]
  and summing these inequalities gives equality:
  \[
    \sum_a\operatorname{Re}\langle x_a,s\rangle
    =
    \operatorname{Re}\langle s,s\rangle
    =
    \|s\|^2
    =
    \|s\|\sum_a\|x_a\|.
  \]
  Hence equality holds in the real-part Cauchy inequality for every \(a\).
  Each \(x_a\) is therefore a non-negative real multiple of \(s\), and the
  multiple is positive because \(x_a\neq0\).  Taking
  \(\xi=s/\|s\|\) proves the claim.
\end{proof}

\begin{lemma}[Equality in the edgewise domination step]
  \label{lem:edgewise-domination-equality}
  Let \(V_\varrho\) carry its unitary Hilbert norm.  Suppose
  \(v=(v_i)_{i\in V}\in\CC^V\otimes V_\varrho\) and
  \(\alpha\in\CC\), \(|\alpha|=\lambda\), satisfy the column equations
  \[
    \alpha v_i=\sum_{e:i\to j}\varrho(\tau(e))v_j
    \qquad(i\in V),
  \]
  and put \(u_i:=\norm{v_i}\).  If the domination argument gives
  \[
    \lambda u_i
    =
    \left\|\sum_{e:i\to j}\varrho(\tau(e))v_j\right\|
    =
    \sum_{e:i\to j}\norm{v_j}
    =
    (Au)_i
    \qquad(i\in V),
  \]
  with all \(u_i>0\), then for each fixed initial vertex \(i\) every
  individual outgoing edge summand
  \[
    x_e:=\varrho(\tau(e))v_{t(e)}
    \qquad(e:i\to t(e))
  \]
  is a positive real multiple of one common vector.  The assertion is
  edgewise: parallel edges and distinct edges with the same terminal vertex
  are separate summands.  More precisely, if
  \(y_i:=\alpha v_i/(\lambda u_i)\), then \(y_i\) is a unit vector and
  \[
    \varrho(\tau(e))v_j=u_j y_i
    \qquad(e:i\to j).
  \]
\end{lemma}

\begin{proof}
  Fix \(i\) and list the outgoing edges with multiplicity as
  \(e_1,\ldots,e_q\).  Put
  \(x_a:=\varrho(\tau(e_a))v_{t(e_a)}\).  Since \(\varrho\) is unitary and
  \(u_{t(e_a)}>0\), each \(x_a\) is non-zero and
  \(\|x_a\|=u_{t(e_a)}\).  The displayed equality is exactly
  \[
    \left\|\sum_{a=1}^q x_a\right\|=\sum_{a=1}^q\|x_a\|.
  \]
  By \Cref{lem:hilbert-triangle-equality}, all of these edge summands are
  positive real multiples of one unit vector.  Since
  \(s:=\sum_a x_a=\alpha v_i\), the common unit vector is \(s/\|s\|\).  Also
  \[
    \sum_{e:i\to j}u_j=(Au)_i=\lambda u_i
    \quad\text{and}\quad
    s=\alpha v_i,
  \]
  so
  \(y_i=s/\|s\|=\alpha v_i/(\lambda u_i)\), and the norm of the \(e:i\to j\)
  summand gives \(x_e=u_j y_i\).  The enumeration was by edge, so the
  conclusion includes parallel edges and repeated terminal vertices.
\end{proof}

\begin{lemma}[Column convention and the determinant block]
  \label{lem:eta-gap-column-determinant-convention}
  Fix an irreducible unitary representation \(\varrho\).  Let
  \(A_\varrho\) denote the row-vector block matrix from
  \eqref{eq:twisted-matrix-def}, which is the matrix used in all trace and
  determinant expressions.  For the line-field argument define
  \[
    \mathcal A_\varrho:=A_\varrho^{\mathrm{col}}
  \]
  as the column block operator on
  \(\CC^V\otimes V_\varrho\) by
  \[
    (\mathcal A_\varrho w)_i
    =
    \sum_{e:i\to j}\varrho(\tau(e))w_j .
  \]
  Then \(\mathcal A_\varrho=A_\varrho^{\mathrm{col}}\) and the determinant block
  \(A_\varrho\) have the same characteristic polynomial.  Consequently
  \[
    \det(tI-\mathcal A_\varrho)
    =
    \det(tI-A_\varrho),
    \qquad
    \rho(\mathcal A_\varrho)=\rho(A_\varrho),
  \]
  and the boundary poles of
  \(\det(I-zA_\varrho)^{-1}\) are exactly the boundary eigenvalues of
  \(\mathcal A_\varrho\) written through this common polynomial.
\end{lemma}

\begin{proof}
  Put
  \(M_{ij}:=\sum_{e:i\to j}\varrho(\tau(e))\).  Choose a basis of
  \(V_\varrho\).  The column operator above has block matrix
  \((M_{ij})_{i,j}\).  The determinant convention treats an element as a
  row family \(x=(x_i)_i\) and multiplies on the right:
  \[
    (xA_\varrho)_j=\sum_i x_iM_{ij}.
  \]
  If the row blocks are vectorized as columns by ordinary transpose, then
  this row action is represented by the full transpose of the column block
  matrix \((M_{ij})_{i,j}\), namely the block matrix with \((j,i)\)-block
  \(M_{ij}^{\mathsf T}\).  Thus no edge is reversed and no label is
  inverted: the same edge \(e:i\to j\) contributes the same linear map
  \(\varrho(\tau(e))\), and only the outside row-versus-column
  realization is transposed.  A finite matrix and its transpose have the
  same characteristic polynomial, so the displayed determinant and
  spectral-radius identities follow.  Substituting \(t=z^{-1}\), or
  equivalently comparing the common polynomial
  \(\det(I-z\,\cdot)\), gives the statement about determinant poles.
\end{proof}

\begin{proposition}[Finite line-field certificate at the Perron boundary]
  \label{thm:eta-gap-rigidity}
  Fix the finite directed graph with primitive adjacency matrix \(A\), the
  fixed finite-group cocycle \(\tau\), and the associated twisted matrices
  \(A_\varrho\) from \eqref{eq:twisted-matrix-def}.  Let
  \(\varrho:G\to U(V_\varrho)\) be an irreducible unitary representation.
  Throughout this proposition the line-field equations use the column block
  convention
  \[
    (\mathcal A_\varrho w)_i
    =
    \sum_{e:i\to j}\varrho(\tau(e))w_j .
  \]
  By \Cref{lem:eta-gap-column-determinant-convention}, this has the same
  characteristic polynomial and spectral radius as the determinant block
  \(A_\varrho\).
  Thus \(\rho(A_\varrho)=\lambda\) if and only if there exist
  \(\omega\in\TT\) and unit vectors
  \(\xi_i\in V_\varrho\), indexed by the vertices \(i\in V\), such that
  \begin{equation}\label{eq:eta-gap-line-field}
    \varrho(\tau(e))\xi_j=\omega\,\xi_i
    \qquad
    \text{for every admissible edge } e:i\to j.
  \end{equation}
  Whenever \eqref{eq:eta-gap-line-field} holds, every closed path
  \(\gamma\) of length \(n\) based at \(i\) satisfies
  \begin{equation}\label{eq:eta-gap-loop-character}
    \varrho(\tau(\gamma))\xi_i=\omega^n\xi_i.
  \end{equation}
  In particular, the subgroup
  \[
    \Gamma_i(\tau)
    :=
    \bigl\langle \tau(\gamma):
      \gamma \text{ closed path based at } i
    \bigr\rangle
  \]
  acts on the line \(\CC\xi_i\) through a one-dimensional character.  This
  proposition is a boundary equality criterion for the fixed finite-state cocycle
  just specified: it identifies exactly when equality
  \(\rho(A_\varrho)=\lambda\) occurs in this fixed non-trivial channel.  Its
  contrapositive is used below only as a finite test for the strict-gap
  specialization in \Cref{cor:strict-gap-specialization}; it is not a classification of cocycles
  for which \(\eta<\lambda\), and it supplies no structural sufficient
  condition on a variable cocycle beyond the stated finite equations.  Thus a
  broad structural theory of sufficient conditions for \(\eta<\lambda\) is not
  imported into the present article.
\end{proposition}

\begin{proof}
  We use \(\mathcal A_\varrho=A_\varrho^{\mathrm{col}}\) throughout the proof. The statement
  has already identified its spectral radius with that of the determinant
  block \(A_\varrho\) by
  \Cref{lem:eta-gap-column-determinant-convention}, so no spectral
  assertion changes under this convention.

  Assume first that \(\rho(A_\varrho)=\lambda\). Since
  \(\mathcal A_\varrho\) is finite dimensional, some eigenvalue
  \(\alpha\in\spec(\mathcal A_\varrho)\) satisfies
  \(|\alpha|=\lambda\). Let
  \(v=(v_i)_{i\in V}\in \CC^V\otimes V_\varrho\) be a corresponding
  eigenvector, so
  \[
    \alpha v_i=(\mathcal A_\varrho v)_i
    =\sum_{e:i\to j}\varrho(\tau(e))v_j
    \qquad (i\in V).
  \]
  Set \(u_i:=\norm{v_i}\). Then
  \[
    \lambda u_i
    =
    \norm{\alpha v_i}
    =
    \left\|\sum_{e:i\to j}\varrho(\tau(e))v_j\right\|
    \le
    \sum_{e:i\to j}\norm{v_j}
    =(Au)_i
  \]
  for every \(i\in V\). Let \(\ell>0\) be a left Perron vector of
  \(A\), normalised by \(\ell^{\mathsf T}A=\lambda\ell^{\mathsf T}\).
  Then
  \[
    \lambda\,\ell^{\mathsf T}u
    \le
    \ell^{\mathsf T}Au
    =
    \lambda\,\ell^{\mathsf T}u,
  \]
  so every coordinate inequality is an equality because \(\ell_i>0\)
  for all \(i\). Hence \(Au=\lambda u\). Primitivity of \(A\) implies
  that \(u\) is a positive Perron eigenvector, so \(u_i>0\) for every
  \(i\in V\).

  Equality in the triangle inequality now holds in every vertex equation.
  Explicitly, for each \(i\),
  \[
    \lambda\norm{v_i}
    =
    \left\|\sum_{e:i\to j}\varrho(\tau(e))v_j\right\|
    =
    \sum_{e:i\to j}\norm{v_j}.
  \]
  Applying \Cref{lem:edgewise-domination-equality} shows that this is an
  equality for the list of outgoing edges with multiplicity, not only after
  grouping by terminal vertex.  Thus no cancellation can occur at a
  Perron-boundary channel, parallel edges are controlled separately, and the
  phase is not chosen separately for each edge; it is the single eigenvalue
  phase \(\omega:=\alpha/\lambda\).  Hence
  \[
    \varrho(\tau(e))v_j
    =
    u_j\,\frac{\alpha v_i}{\lambda u_i}
    \qquad
    \text{for every admissible edge } e:i\to j.
  \]
  Write \(\xi_i:=v_i/u_i\in V_\varrho\). Dividing by \(u_j\) yields
  \[
    \varrho(\tau(e))\xi_j=\omega\xi_i,
    \qquad e:i\to j,
  \]
  with unit vectors \(\xi_i\).  This gives
  \eqref{eq:eta-gap-line-field}.  Thus the equality case is not only a
  norm-level statement: equality in every vertex domination inequality forces
  equality in every edgewise Hilbert-space triangle inequality, and hence the
  same phase \(\omega=\alpha/\lambda\) relates all outgoing edge summands.

  Conversely, assume that \eqref{eq:eta-gap-line-field} holds. Let
  \(h=(h_i)_{i\in V}>0\) be a right Perron vector of \(A\), so
  \(Ah=\lambda h\). Define \(w=(w_i)_{i\in V}\) by
  \(w_i:=h_i\xi_i\). Then
  \begin{align*}
    (\mathcal A_\varrho w)_i
    &=
    \sum_{e:i\to j}\varrho(\tau(e))w_j \\
    &=
    \sum_{e:i\to j}h_j\varrho(\tau(e))\xi_j \\
    &=
    \omega\sum_{e:i\to j}h_j\xi_i \\
    &=
    \omega(Ah)_i\xi_i \\
    &=
    \omega\lambda w_i.
  \end{align*}
  Thus \(\omega\lambda\in\spec(\mathcal A_\varrho)\), so
  \(\rho(\mathcal A_\varrho)\ge \lambda\), equivalently
  \(\rho(A_\varrho)\ge \lambda\). On the other hand,
  \Cref{prop:twisted-spectral-radius-bound} gives
  \(\rho(A_\varrho)\le \lambda\). Hence
  \(\rho(A_\varrho)=\lambda\).

  Finally, iterating \eqref{eq:eta-gap-line-field} along a closed path
  \(\gamma=e_1\cdots e_n\) based at \(i\) gives
  \eqref{eq:eta-gap-loop-character}. Hence the line \(\CC\xi_i\) is
  \(\Gamma_i(\tau)\)-invariant. The action of \(\Gamma_i(\tau)\) on
  this line defines a one-dimensional character.
\end{proof}

\begin{corollary}[Finite-state diagnostic for the strict twisted gap]
  \label{cor:eta-gap-boundary-classifier}
  The strict-gap hypothesis used by \Cref{cor:strict-gap-specialization} and the clean computational formulas is equivalent, for
  the fixed finite extension, to the absence of the following Perron-boundary
  line fields.
  Suppose \(G\) is non-trivial. Then the following are equivalent.
  \begin{enumerate}
    \item The ordinary strict-gap Perron-point endpoint chart fails,
      namely \(\eta=\lambda\).
    \item There exist a non-trivial irreducible representation
      \(\varrho\in\Irr(G)\), a phase \(\omega\in\TT\), and unit vectors
      \(\xi_i\in V_\varrho\), one for each vertex, such that
      \begin{equation}\label{eq:eta-gap-boundary-classifier}
        \varrho(\tau(e))\xi_j=\omega\xi_i
        \qquad
        \text{for every admissible edge } e:i\to j.
      \end{equation}
  \end{enumerate}
  Consequently, for this fixed finite-state cocycle, \(\eta<\lambda\) holds
  exactly when the finite list of edge equations
  \eqref{eq:eta-gap-boundary-classifier} has no solution in any non-trivial
  irreducible representation. In this case the strict-gap hypotheses of
  \Cref{thm:class-euler-product-mertens} are reduced to the verified absence
  of these boundary line fields.  If the equations do have a solution, the
  primary Perron-boundary package remains governed instead by
  \Cref{def:peripheral-peter-weyl-data,lem:finite-cover-peripheral-expansion,thm:renormalized-frobenius-product}.
\end{corollary}

\begin{proof}
  By \Cref{prop:twisted-spectral-radius-bound}, every non-trivial channel
  satisfies \(\rho(A_\varrho)\le\lambda\). Since \(G\) is finite, there
  are only finitely many irreducible representations, so the maximum
  defining \(\eta\) is achieved. Hence \(\eta=\lambda\) if and only if
  \(\rho(A_\varrho)=\lambda\) for at least one non-trivial
  \(\varrho\in\Irr(G)\). Applying \Cref{thm:eta-gap-rigidity} to that
  representation gives exactly the edge equations
  \eqref{eq:eta-gap-boundary-classifier}, and the converse implication in
  the same theorem shows that any solution of those equations forces
  \(\rho(A_\varrho)=\lambda\), hence \(\eta=\lambda\). The final sentence
  is the contrapositive, together with the role of the strict-gap
  hypothesis stated in \Cref{thm:class-euler-product-mertens}.  It says
  nothing against the renormalized Perron-boundary finite-part theorem,
  which is invoked through its own peripheral certificate.
\end{proof}

\begin{proposition}[Finite-state diagnostic certificate for the Perron-point
  hypothesis]
  \label{prop:eta-gap-finite-state-certificate}
  Assume \(\lambda>1\). Let \(\widetilde A\) be the skew-product
  adjacency matrix attached to \(\tau\), and set
  \[
    \chi_{\widetilde A}(t):=\det(tI-\widetilde A).
  \]
  The following finite-state conditions are equivalent.  Thus the strict-gap
  specialization may be checked for the fixed cocycle by asking for either the
  spectral-radius condition in \textup{(ii)}/\textup{(iv)} or the line-field
  exclusion in \textup{(iii)}, rather than treating \(\eta<\lambda\) as an
  unverifiable black-box hypothesis.
  \begin{enumerate}[label=\rm(\roman*)]
  \item The strict twisted spectral gap used in the ordinary endpoint
    class-Mertens and Frobenius-class product formulas holds:
    \[
      \eta=\max_{\varrho\neq\mathbf 1}\rho(A_\varrho)<\lambda .
    \]
  \item The skew-product characteristic polynomial has exactly one root on
    the Perron circle, namely a simple root at \(\lambda\); equivalently
    \(\chi_{\widetilde A}(t)=(t-\lambda)q(t)\) with every root of the
    finite polynomial \(q\) having modulus strictly less than \(\lambda\).
  \item No non-trivial irreducible representation admits a boundary line
    field: for every \(\varrho\neq\mathbf1\), the finite algebraic system
    \[
      \varrho(\tau(e))\xi_j=\omega\xi_i\quad(e:i\to j),
      \qquad \omega\in\TT,
      \qquad \norm{\xi_i}=1\quad(i\in V)
    \]
    has no solution.
  \item The Perron-circle certificate returns \(\mathsf{Pass}\).  This
    certificate is stated first in the integral skew-product model and is not
    dependent on choosing unitary irreducible representation matrices.  It
    computes the integer polynomial
    \(\chi_{\widetilde A}(t)\in\ZZ[t]\), removes the distinguished factor
    \(t-\lambda\), isolates the finitely many remaining algebraic roots in
    disjoint rational boxes in the splitting field of
    \(\chi_{\widetilde A}\), and verifies by exact algebraic comparison that
    every remaining root has modulus strictly smaller than \(\lambda\).  If
    one of these comparisons fails, the certificate returns
    \(\mathsf{Fail}\).  A block-level audit may then locate the failed root
    inside a non-trivial Peter--Weyl block, but this localization is not part
    of the primary pass certificate and must use the exact block model of
    \Cref{lem:eta-gap-exact-block-models}.
  \end{enumerate}
  Consequently the Perron-point assumptions in
  \Cref{thm:class-mertens-explicit,thm:class-euler-product-mertens} are
  exactly \(\lambda>1\) plus any one of the finite certificates above.
  The pass/fail output is a diagnostic for this finite datum: a pass certifies the simultaneous
  availability of all non-trivial Perron-point determinant branches, while a
  fail is equivalent to the existence of a non-trivial boundary line field and
  hence to a Perron-circle pole in some non-trivial determinant factor.
  Under these certificates the constants
  \(\mathcal M_C(A,\tau;G)\) and \(K_C\) are those supplied by
  \Cref{thm:class-mertens-explicit,thm:class-euler-product-mertens}; if the
  equivalent conditions fail, this paper asserts no ordinary unrenormalized
  Perron-point evaluation for the non-trivial primitive channels in the
  strict-gap chart.  Certified boundary collisions are handled instead by the
  reduced finite parts \(R_\pi\), \(F_\pi^{\mathrm{ren}}\), \(\delta_C\), and
  \(K_C^{\mathrm{ren}}\) of
  \Cref{thm:perron-boundary-determinant-product}.
\end{proposition}

\begin{proof}
  The equivalence of \rm(i) and \rm(ii) is exactly the skew-product
  peripheral-spectrum criterion of
  \Cref{lem:twisted-gap-skew-product-peripheral}: the Peter--Weyl block
  decomposition of \(\widetilde A\) gives the spectrum of
  \(\widetilde A\), with algebraic multiplicities, as the disjoint union
  of the spectra of the blocks \(A_\varrho\), each repeated
  \(\dim\varrho\) times. The trivial block is \(A\), and primitivity of
  \(A\) makes \(\lambda\) the only peripheral eigenvalue of that block and
  makes it simple. Thus all remaining roots of
  \(\chi_{\widetilde A}\) lie inside \(|t|<\lambda\) precisely when every
  non-trivial block has spectral radius strictly smaller than \(\lambda\).

  The equivalence of \rm(i) and \rm(iii) is
  \Cref{cor:eta-gap-boundary-classifier} when non-trivial irreducible
  representations exist. If none exist, then \(G\) is trivial, \(\eta=0\)
  by convention, and the finite line-field condition is vacuous, so both
  \rm(i) and \rm(iii) hold because \(\lambda>1\). In the remaining case,
  \Cref{thm:eta-gap-rigidity} says that equality
  \(\rho(A_\varrho)=\lambda\) for a fixed non-trivial irreducible
  representation is equivalent to the displayed line-field equations for
  that representation. Since \(G\) is finite, only finitely many
  irreducible representations occur, and the maximum defining \(\eta\) is
  attained. Therefore the strict gap fails exactly when one of these
  finite systems has a solution.

  Item~\rm(iv) is the same criterion written as a finite certificate in the
  integer skew-product model.  The matrix \(\widetilde A\) has entries in
  \(\ZZ_{\ge0}\), hence \(\chi_{\widetilde A}\in\ZZ[t]\).  The Perron root
  \(\lambda\) is the simple Perron root of the base primitive matrix \(A\);
  exact comparison may be made in the number field
  \(\QQ(\lambda,\alpha)\), where \(\alpha\) is the root being tested, or
  equivalently by a square-free factor/root-isolation record for
  \(\chi_{\widetilde A}(t)/(t-\lambda)\).  Such a record consists of
  isolating rational boxes, their multiplicities, and exact sign
  comparisons for
  \[
    \lambda^2-\alpha\overline{\alpha}>0 .
  \]
  This is a finite algebraic comparison; no numerical choice of a
  representation model enters.  Therefore the certificate returns
  \(\mathsf{Pass}\) exactly when item~\rm(ii) holds.  If it returns
  \(\mathsf{Fail}\), item~\rm(ii) fails; by the already proved equivalence
  with \rm(i) and \rm(iii), some non-trivial channel has a boundary line
  field.  If one wants to identify the failed channel explicitly, one may
  refine the certificate by computing the Peter--Weyl blocks over a
  cyclotomic field containing the representation matrices, factoring the
  block characteristic polynomials there, and checking the same root
  inclusion and rank/eigenspace tests exactly under a specified complex
  embedding.  This optional refinement is
  a localization of the integral skew-product certificate, not a replacement
  for it.

  The final assertion is a substitution of this certificate into the
  hypotheses of
  \Cref{thm:class-mertens-explicit,thm:class-euler-product-mertens}. Those
  theorems require \(\lambda>1\) and \(\eta<\lambda\) for the Perron-point
  evaluation of the non-trivial primitive Peter--Weyl channels; the present
  proposition only identifies that hypothesis with finite matrix or
  finite algebraic tests.  None of these boundary criteria supplies an
  additional constant formula when the tests fail, and no general
  classification of cocycles forcing \(\eta<\lambda\) is used below.  When
  the tests fail, the relevant replacement is not another strict endpoint
  value but the Perron-boundary renormalization package with its finite
  algebraic multiplicity and branch data.
\end{proof}

\begin{lemma}[Exact block models for optional localization]
  \label{lem:eta-gap-exact-block-models}
  The strict-gap certificate of
  \Cref{prop:eta-gap-finite-state-certificate} is independent of any chosen
  matrix model for the irreducible representations.  If a computation uses
  Peter--Weyl blocks rather than the integer skew-product matrix, the exact
  algebraic model is fixed as follows.  Let \(e_G\) be the exponent of \(G\).
  Every finite-group representation may be realized over the cyclotomic field
  \(K=\QQ(\zeta_{e_G})\).  Choose matrices
  \(\pi(g)\in\operatorname{GL}_{d_\pi}(K)\) for all \(g\in G\), form
  \(A_\pi\) over \(K\), and compare eigenvalues after a specified embedding
  \(K\hookrightarrow\CC\).  Peripheral membership, algebraic multiplicity,
  and semisimplicity are then finite exact tests:
  factor \(\det(tI-A_\pi)\) over \(K\), isolate the roots of each factor under
  the chosen embedding, compare \(|\alpha|\) with \(\lambda\) in the compositum
  \(K(\lambda,\alpha,\overline\alpha)\), and, when needed, compute
  \(\dim_K\ker(A_\pi-\alpha I)\) by exact rank over \(K(\alpha)\).  Equivalent
  nonunitary matrices may be used for these determinant and rank tests, while
  the line-field domination argument itself is read in a unitarized equivalent
  Hilbert model.
\end{lemma}

\begin{proof}
  The field-of-definition statement is the standard finite-group fact that
  the values and matrix entries of a complete set of irreducible
  representations may be taken in a cyclotomic field of exponent dividing
  \(e_G\).  Passing from one equivalent model to another conjugates every
  block \(A_\pi\), so characteristic polynomials, algebraic multiplicities,
  eigenspace dimensions, and rank tests are unchanged after extension of
  scalars.  Exact root isolation and the sign of
  \(\lambda^2-\alpha\overline\alpha\) are algebraic questions in the displayed
  compositum.  Finally, because \(G\) is finite, every complex representation
  is equivalent to a unitary one by averaging an inner product over \(G\).
  Thus nonunitary equivalent matrices are legitimate for exact determinant
  computations, whereas the norm-equality proof in
  \Cref{thm:eta-gap-rigidity} uses the corresponding unitary model.
\end{proof}

\begin{proposition}[Perron-point separation for the strict-gap specialization]
  \label{prop:eta-gap-perron-separation}
  Assume \(\lambda>1\).  The strict twisted spectral gap
  \(\eta<\lambda\) is equivalent to each of the following Perron-point
  analytic conditions.
  \begin{enumerate}[label=\textup{(\roman*)},leftmargin=*]
  \item For every non-trivial irreducible representation \(\varrho\), the
    branch
    \[
      \log\det(I-zA_\varrho)^{-1}
      =
      \sum_{n\ge1}\frac{\Tr(A_\varrho^n)z^n}{n}
    \]
    used in the primitive determinant theorem extends holomorphically to a neighbourhood of
    the closed disc \(|z|\le\lambda^{-1}\).
  \item For every non-trivial class function \(\phi\) with zero average on
    \(G\), all non-trivial Peter--Weyl determinant coordinates entering the
    Adams--M\"obius formula for \(\Pi_\phi\) have branch-compatible values at
    \(z=\lambda^{-1}\), and hence the Abel-path Perron values used in
    \Cref{thm:class-mertens-explicit,thm:class-euler-product-mertens}, equivalently the
    strict-gap specialization \Cref{cor:strict-gap-specialization}, are
    computed from the branch-compatible determinant polynomials and the
    recorded strict-gap branch convention.
  \end{enumerate}
  If the strict gap fails, then there are a non-trivial irreducible
  representation \(\varrho\), a phase \(\omega\in\TT\), and a boundary line
  field satisfying \eqref{eq:eta-gap-line-field}.  Moreover
  \(\omega\lambda\) is a zero of the common characteristic polynomial from
  \Cref{lem:eta-gap-column-determinant-convention}, so the non-trivial
  determinant factor \(\det(I-zA_\varrho)^{-1}\) has a pole at
  \(z=(\omega\lambda)^{-1}\) on the Perron circle.  Thus the finite-state
  certificate of \Cref{prop:eta-gap-finite-state-certificate} is a necessary
  and sufficient verification gate for the ordinary Perron-point constants in
  the strict-gap specialization; it is not an automatic existence theorem for
  arbitrary cocycles and it is not a gate for the primary renormalized
  boundary constants.
\end{proposition}

\begin{proof}
  The equivalence between \(\eta<\lambda\) and item~\textup{(i)} is purely
  finite dimensional.  If \(\eta<\lambda\), then every eigenvalue \(\alpha\)
  of every non-trivial block \(A_\varrho\) satisfies
  \(|\alpha|\le\rho(A_\varrho)<\lambda\).  Hence
  \(1/\alpha\), for \(\alpha\neq0\), lies strictly outside the closed disc
  \(|z|\le\lambda^{-1}\).  The determinant
  \(\det(I-zA_\varrho)
  =\prod_\alpha(1-z\alpha)\) therefore has no zero on a neighbourhood of
  that closed disc, and the logarithm branch normalized at \(0\) extends
  holomorphically there.

  Conversely, if item~\textup{(i)} holds and some non-trivial
  \(A_\varrho\) had \(\rho(A_\varrho)=\lambda\), then an eigenvalue
  \(\alpha\in\spec(A_\varrho)\) would satisfy \(|\alpha|=\lambda\).  The
  determinant \(\det(I-zA_\varrho)\) would vanish at
  \(z=\alpha^{-1}\), a point of the Perron circle, so the
  reciprocal determinant and its normalized logarithm could not be
  holomorphic in any neighbourhood of the closed Perron disc.  This
  contradicts item~\textup{(i)}.  Since
  \Cref{prop:twisted-spectral-radius-bound} gives
  \(\rho(A_\varrho)\le\lambda\) for every \(\varrho\), item~\textup{(i)}
  forces \(\eta<\lambda\).

  Item~\textup{(ii)} is exactly the same condition after finite character
  decomposition.  The zero-average class functions have no trivial
  character coordinate.  Their periodic logarithms are finite linear
  combinations of the non-trivial determinant logarithms by
  \Cref{thm:class-function-linearisation}, and the primitive series are then
  obtained from those finitely many logarithms by the Adams--M\"obius formula
  of \Cref{cor:class-function-adams-mobius}.  Therefore the Perron values
  used by \Cref{thm:class-mertens-explicit,thm:class-euler-product-mertens}
  are available from the branch-compatible determinant data precisely when the
  non-trivial determinant branches in item~\textup{(i)} extend through the
  closed Perron disc.  The trivial channel is deliberately excluded here;
  it is singular at \(z=\lambda^{-1}\) and is handled separately by the
  scalar finite-part normalisation.

  Finally suppose the strict gap fails.  Because only finitely many
  irreducibles occur and each non-trivial block has spectral radius at most
  \(\lambda\), there is a non-trivial \(\varrho\) with
  \(\rho(A_\varrho)=\lambda\).  \Cref{thm:eta-gap-rigidity} supplies a phase
  \(\omega\) and a boundary line field satisfying
  \eqref{eq:eta-gap-line-field}; the converse half of its proof constructs
  an eigenvector of the column realization with eigenvalue
  \(\omega\lambda\).  By
  \Cref{lem:eta-gap-column-determinant-convention}, this is a zero of the
  determinant block's characteristic polynomial.  Hence
  \(\det(I-zA_\varrho)=0\) at \(z=(\omega\lambda)^{-1}\), and the reciprocal
  determinant has a boundary pole.  This is precisely the obstruction
  detected by \Cref{prop:eta-gap-finite-state-certificate}.
\end{proof}

\paragraph{Interface with the strict-gap specialization.}
Combining
\Cref{cor:eta-gap-boundary-classifier,prop:eta-gap-finite-state-certificate,prop:eta-gap-perron-separation}
gives the only export from this section.  For the fixed finite cocycle, the
strict-gap chart used in
\Cref{cor:strict-gap-specialization,thm:class-mertens-explicit,thm:class-euler-product-mertens}
is available exactly when the finite certificate eliminates every non-trivial
boundary line field.  On that branch the non-trivial determinant logarithms
have the ordinary Perron-point values used by the strict-gap formulas.  If the
certificate fails, some non-trivial determinant factor has a pole on the
Perron circle; the ordinary endpoint chart is then unavailable, while the main
Perron-boundary product theorem uses the certified peripheral data and reduced
finite parts instead.

\begin{remark}[Strict-gap diagnostic boundary]
  \label{rem:eta-gap-criteria-scope}
  \label{rem:eta-gap-criteria-boundary-note}
  The diagnostics above characterize, for a fixed finite cocycle, the boundary
  \(\eta=\lambda\) by a finite line-field certificate and provide the hypothesis
  check needed for \Cref{cor:strict-gap-specialization} and the older strict-gap
  formulas.
\end{remark}


